%% fig-planes.tex — Backup R2: Two-planes revocation diagram.
%% Answers Q4 (stateless JWT) — shown on demand during Q&A.
%% B1 fixes applied:
%%   - English «terminated» → «розірвано»; «cleanup» → «очищення» (all Ukrainian).
%%   - 60с note moved to y=-2.35 (was -1.9 = box bottom ≈ struck by border).
%%   - Plane titles lifted to y=2.15 (were 1.7 = inside box top border).
%%   - planeBox height enlarged to 4.4cm so internal steps have more air.
%% Assumes styles.tex pre-loaded.
\begin{tikzpicture}[x=1cm, y=1cm,
  planeBox/.style={%
    draw=clInk, fill=clBg, line width=1.2pt, rounded corners=4pt,%
    minimum width=5.8cm, minimum height=4.4cm, align=center, inner sep=8pt%
  },%
  stepNode/.style={%
    draw=clInk, fill=clBg, line width=1pt, rounded corners=2pt,%
    font=\small\sffamily, align=center, inner sep=4pt,%
    minimum width=4.6cm%
  }%
]
%%
%% ── LEFT PLANE: Control plane ────────────────────────────────────────────────
\node[planeBox] (CP) at (0, 0) {};
%% Title lifted above box top (box half-height = 2.2; title at 2.15 = just above) ─
\node[font=\small\sffamily\bfseries, text=clInk, above] at (0, 2.15)
  {Площина керування};
%%
%% Steps inside
\node[stepNode] (jwt) at (0, 0.85) {JWT stateless\\(без blacklist)};
\node[stepNode, draw=clAccent, fill=clAccent!8] (rev) at (0, -0.3)
  {API: відкликати сесію};
%% Shrink font + shorten text so node sits fully inside planeBox (5.8cm wide)
\node[stepNode, draw=clAccent, fill=clAccent!8,
      font=\scriptsize\ttfamily] (del) at (0, -1.3)
  {wg set \ldots{} remove};
%%
\draw[->, >=Stealth, line width=1pt, clInk] (jwt.south) -- (rev.north);
\draw[->, >=Stealth, line width=1pt, clAccent] (rev.south) -- (del.north);
%%
%% 60с note — moved clear below box bottom (was at -1.9 = box border; fix B1) ──
\node[font=\scriptsize\sffamily, text=clGrey, align=center] at (0, -2.35)
  {(або: фонове очищення кожні 60\,с)};
%%
%% ── RIGHT PLANE: Data plane ───────────────────────────────────────────────────
\node[planeBox] (DP) at (7.2, 0) {};
%% Title lifted above box top ─────────────────────────────────────────────────
\node[font=\small\sffamily\bfseries, text=clInk, above] at (7.2, 2.15)
  {Площина даних};
%%
%% wg0 bus inside data plane
\draw[wgbus] (4.8, 0.85) -- (9.6, 0.85);
\node[font=\footnotesize\ttfamily, text=clGrey, right] at (9.65, 0.85)
  {\texttt{wg0}};
%%
%% Present peer (filled accent)
\node[wgpeer, minimum width=2.6cm, font=\scriptsize\sffamily\bfseries]
  (peerOn) at (6.3, 0.85) {peer\,\cmark};
%%
%% Gone peer (hollow grey)
\node[wggone, minimum width=2.6cm, font=\scriptsize\sffamily, text=clGrey]
  (peerOff) at (8.1, 0.85) {peer};
%%
%% Remove arrow between states
\draw[wgremove] (peerOn.east) -- (peerOff.west);
\node[font=\scriptsize\sffamily, text=clGrey, above] at (7.2, 1.0)
  {remove};
%%
%% Effect labels — Ukrainian only (fix B1) ──────────────────────────────────
\node[font=\small\sffamily\bfseries, text=clAccent, align=center]
  at (7.2, -0.25) {Тунель розірвано\\(миттєво)};
%%
\node[font=\scriptsize\sffamily, text=clGrey, align=center]
  at (7.2, -1.05) {навіть з валідним JWT:\\трафік обрізано на рівні ядра};
%%
%% ── Arrow between planes ──────────────────────────────────────────────────────
\draw[->, >=Stealth, line width=1.5pt, clAccent]
  ($(del.east)+(0.08,0)$) -- ($(DP.west)+(0,-1.3)$);
\node[font=\scriptsize\sffamily\bfseries, text=clAccent, above]
  at (3.65, -1.3) {синхронно};
%%
%% ── Figure title ──────────────────────────────────────────────────────────────
\node[font=\normalsize\sffamily\bfseries, text=clInk, align=center]
  at (3.6, 2.9) {Відкликання: дві площини};
%%
\end{tikzpicture}
